\problem{1}
\textbf{[Отбор Сириус 2024, 7.4]}
В селе Весёлое проживают 3700 жителей. Однажды известный в селе математик заметил, что:

- для каждого человека, проживающего в селе, можно найти его ровесника (человека такого же возраста) в этом селе;

- в любой группе из 192 жителей села найдутся трое одного возраста.

Докажите, что на местную дискотеку можно пригласить 39 жителей села так, чтобы все они были одного возраста.

\problem{2}
\textbf{[Отбор Сириус 2024, 7.6]}
В чемпионате по футболу участвуют 30 команд. Каждая команда состоит либо только из рыцарей, которые всегда говорят правду, либо только из лжецов, которые всегда лгут. В первом туре команды были разбиты на 15 пар и каждая пара сыграла матч. На пресс-конференции после тура капитаны десяти команд заявили, что их команды проиграли, капитаны других десяти команд - что их команды выиграли, а капитаны оставшихся десяти команд - что их команды сыграли вничью. Известно, что в первом туре победы могли одержать лишь команды рыцарей. Какое наибольшее количество команд-лжецов могло участвовать в чемпионате?

\problem{3}
\textbf{[Отбор Сириус 2021, 7.5]}
Незнайка очень обеспокоен своей средней оценкой по математике. Она оказалась меньше 3. (Каждая оценка по математике - целое число от 1 до 5 ). Для улучшения оценок он отработал несколько контрольных работ и исправил оценки по ним с «единицы» на «тройку». Докажите, что средняя оценка Незнайки не достигла «четверки».

\problem{4}
\textbf{[Отбор Сириус 2021, 7.6]}
В таблице с 4 строками и $k$ столбцами записаны натуральные числа так, что числа в каждой строке различные и сумма чисел в каждом столбце равна 30 . Найдите максимально возможное количество столбцов в таблице.

\problem{5}
\textbf{[Отбор Сириус 2023, 7.1]}
Вершины правильного многоугольника занумерованы натуральными числами начиная с 1 в порядке возрастания по часовой стрелке. При этом 7 -ая вершина оказалась напротив 23-ей. Сколько всего вершин в этом многоугольнике?

\newpage

\hspace{3mm}

\problem{6}
\textbf{[Отбор Сириус 2022, 7.6]}
В одной из профильных смен Сириуса обучались 123 школьника. У каждого из них есть хотя бы один друг среди других ребят на этой смене. Покажите, что для любого четного числа $1<n<123$ в этой смене можно выбрать такую группу из $n$ учеников, что у каждого входящего в группу ученика нашёлся хотя бы один друг в этой же группе.